\section{Discussão e limitações}
\label{sec:discussao}

\paragraph{Interpretação.}
Os resultados da Seção~\ref{sec:resultados} não devem ser lidos como ``verdades'' sobre equipes reais: mostram \textbf{comportamento sensível} às regras explícitas do motor sob um desenho experimental mínimo. Recorde-se a distinção da Seção~\ref{sec:sinergia-neutral-vs}: relativamente à visão com $S\equiv 0$, introduzir sinergias não ``muda o jogo'' de propósito (continuam FIFO, WIP e sorteios de capacidade), mas \textbf{relocaliza} parte da variabilidade de fluxo em pares explícitos e nos mecanismos de colaboração/repasse. Em particular, a sinergia negativa de repasse (cenário C) combina o multiplicador de repasse com maior propensão a retrabalho, enquanto a sinergia positiva no par colaborativo (cenário B) amplifica o trabalho efetivo apenas na coluna \texttt{dev} para cartões colaborativos. O cenário D ilustra a combinação dos dois canais de $S$, reforçando a necessidade de desenhos experimentais mais amplos quando o objetivo é estimar magnitude com precisão.

\paragraph{Reprodutibilidade e transparência.}
Manter o motor em TypeScript puro e versionar em conjunto (i) parâmetros, (ii) cenários, (iii) scripts de exportação do artigo e (iv) a aplicação web que consome o mesmo pacote \texttt{src/simulation} reduz ambiguidade sobre ``o que foi simulado'', alinhando-se às práticas recomendadas por Wilson et al.~\cite{wilson2017good} e às regras de Sandve et al.~\cite{sandve2013ten} para pesquisa computacional reprodutível. Trabalhos futuros devem publicar \emph{releases} com \texttt{package-lock.json} (ou equivalente) e um identificador de \emph{commit} ao reproduzir números em artigos ou materiais didáticos.

\paragraph{Limitações.}
O modelo é redutivo: sem aprendizado, sem filas paralelas de defeitos, sem mudanças de prioridade no meio da sprint, e a retrospectiva não altera parâmetros. A sinergia $s_{ij}$ é um escalar-resumo por par.

\paragraph{Validação.}
Não reportamos validação empírica com dados de projetos reais; esse é um passo natural seguinte (calibração de $\beta$ e $\gamma$), como discutido por Kellner et al.~\cite{kellner1999software} no contexto de simulação de processos de software.
